% =============================================================================
% Chapter 9: Conclusion, Limitations, and Future Work
% Pit-Viper-Inspired Infrared-Thermal and Visual Sensor Fusion
% for Nocturnal UAV Navigation
% =============================================================================
\selectlanguage{hebrew}

\section{סיכום, מגבלות ועבודה עתידית}
\label{sec:ch09_conclusion}

\subsection{סיכום התרומות המרכזיות}
\label{sec:ch09_contributions}

במאמר זה הצגנו מערכת מיזוג תרמי-חזותי מלאה לניווט רחפנים לילי, בהשראת מנגנון קליטת האינפרה-אדום של הצפע הגומתי. המערכת, הקרויה \en{ThermalVisionFusion} (\en{TVF}), משלבת עקרונות ביולוגיים עם טכניקות עיבוד תמונה ולמידה עמוקה מתקדמות, תוך התאמה לדרישות ההפעלה בזמן אמת על פלטפורמת רחפן מוגבלת משאבים.

התרומה המדעית העיקרית של מחקר זה היא פיתוח מנגנון מיזוג אדפטיבי מבוסס רשת קשב צולב (\en{cross-attention}) המתאים את משקל המיזוג $\alpha(t)$ בהתאם לתנאי התאורה הסביבתיים. מנגנון זה מאפשר ניצול אופטימלי של המידע משני התחומים הספקטרליים: בתנאי חושך מוחלט (מתחת ל-0.1 \en{lux}), המשקל התרמי גבוה ומאפשר זיהוי מכשולים על בסיס חתימת החום שלהם; בתנאי תאורה טובים (מעל 10 \en{lux}), המשקל החזותי גובר ומאפשר ניצול המרקם והצבע לזיהוי מדויק יותר.

תוצאות הניסוי המפורטות בפרק 8 מדגימות את יעילותה של המערכת. \en{ATE} הממוצע של מערכת \en{TVF} הוא 1.15 מטרים, שיפור של 73\% לעומת מערכת \en{ORB-SLAM3} חזותית בלבד (4.28 מטרים) ו-60\% לעומת \en{LOAM} (2.87 מטרים). שיעור אבדן העקיבה ירד מ-38.5\% ל-3.1\% בלבד, ושיעור ההימנעות ממכשולים הגיע ל-94.7\%. כל ההבדלים מובהקים סטטיסטית ברמת $p < 0.001$.

תרומה נוספת היא פיתוח שיטת כיול משותפת למצלמה חזותית ותרמית, המאפשרת רישום גיאומטרי מדויק עם שגיאה של 0.8 פיקסלים בלבד. שיטת הכיול מבוססת על לוח ייעודי הנראה בשני התחומים הספקטרליים, ומשלבת כיול אינטרינזי ואקסטרינזי בתהליך אחיד. כמו כן, פותח אלגוריתם סנכרון זמני מבוסס קורלציה צולבת של אותות תנועה, המשיג דיוק של 2.1 מילישניות.

מבחינה מעשית, המערכת מומשה על פלטפורמת \en{NVIDIA Jetson Orin NX} עם צריכת הספק ממוצעת של 22.3 ואט וזמן השהיה כולל של 26.8 מילישניות, העומד בדרישות קצב העדכון של 30 \en{Hz}. המערכת הודגמה בהצלחה על 12 משימות טיסה ליליות בשטח פתוח, תוך הפגנת עמידות בתנאי תאורה משתנים ובנוכחות מכשולים מסוגים שונים.

\subsection{מגבלות המערכת הנוכחית}
\label{sec:ch09_limitations}

למרות הביצועים המעולים, זוהו מספר מגבלות משמעותיות במערכת הנוכחית. ראשית, המערכת תלויה בכיול מדויק בין המצלמות, העלול להשתנות במהלך טיסה ממושכת עקב רעידות, שינויי טמפרטורה, או פגיעה מכנית. הכיול הנוכחי דורש כיול מחדש כל 10 שעות טיסה בממוצע, מה שמגביל את האוטונומיה של המערכת במשימות ארוכות טווח.

שנית, ביצועי המערכת יורדים בתנאי מזג אוויר קיצוניים. גשם כבד, ערפל, או אבק צפוף מחלישים את האות התרמי ומפחיתים את טווח הזיהוי האפקטיבי מ-150 מטרים ל-40 מטרים בלבד. כמו כן, טמפרטורת סביבה גבוהה (מעל 40$^{\circ}$C) מפחיתה את הניגודיות התרמית בין גופים חמים לסביבה, ומקשה על זיהוי מכשולים.

מגבלה שלישית נוגעת ליכולת ההכללה (\en{generalization}) של מודל רשת המיזוג. המודל אומן על נתונים שנאספו בנגב המערבי, והביצועים שלו יורדים ב-12-18\% בסביבות שונות מבחינה גיאוגרפית ואקלימית. איסוף נתונים ממגוון רחב יותר של סביבות גיאוגרפיות ישפר את יכולת ההכללה של המודל.

מגבלה נוספת היא התלות במצלמה תרמית ברזולוציה גבוהה יחסית (640x512), המוסיפה משקל של 180 גרם ועלות של כ-5,000 דולר למערכת. עבור רחפנים קטנים במיוחד (\en{nano-drones}), משקל ועלות אלה עשויים להיות בלתי סבירים. פיתוח גרסה קלת משקל של המערכת המתבססת על מצלמה תרמית ברזולוציה נמוכה יותר (160x120) היא כיוון מחקר עתידי חשוב.

לבסוף, המערכת הנוכחית אינה מנצלת מידע סמנטי אודות הסביבה. זיהוי מכשולים מתבצע על בסיס חתימה תרמית וגיאומטרית בלבד, ללא הבנה של קטגוריית המכשול או התנהגותו הצפויה. שילוב של מודל זיהוי סמנטי (\en{semantic segmentation}) עשוי לשפר את קבלת ההחלטות במצבים מורכבים.

\subsection{כיווני מחקר עתידיים}
\label{sec:ch09_future_work}

בהתבסס על המגבלות שזוהו ועל המגמות העדכניות בתחום, אנו מציעים מספר כיווני מחקר עתידיים מבטיחים. ראשית, שילוב של חיישני אירועים (\en{event-based cameras}) לצד המצלמה התרמית עשוי לשפר משמעותית את ביצועי המערכת בתנאי תאורה משתנים במהירות. חיישני אירועים, המגיבים לשינויים בעוצמת האור בזמן אמת עם השהיה של מיקרושניות בודדות, יכולים לספק מידע משלים למידע התרמי במקרים של מעבר מהיר בין אזורים מוארים לחשוכים.

כיוון מחקר שני הוא שילוב של ראייה קוונטית (\en{quantum imaging}) לזיהוי מכשולים בתנאי חושך מוחלט. טכנולוגיות ראייה קוונטית, המבוססות על גילוי פוטונים בודדים, יכולות לפעול בתנאי תאורה נמוכים בהרבה מאלה הנדרשים למצלמה תרמית קונבנציונלית. עם זאת, טכנולוגיה זו נמצאת עדיין בשלבי מחקר מוקדמים ודורשת פיתוח נוסף לפני שניתן יהיה לשלבה במערכת רחפן.

כיוון מחקר שלישי הוא הרחבת המערכת לניווט נחילי רחפנים (\en{swarm navigation}). נחיל רחפנים המצויד במערכות מיזוג תרמי-חזותי יכול לכסות שטח נרחב יותר, לספק ראייה רב-כיוונית, ולשתף מידע על מכשולים בזמן אמת. האתגר המרכזי בכיוון זה הוא פיתוח פרוטוקולי תקשורת יעילים ומודלים לשיתוף מידע בין רחפנים תוך שמירה על עקביות המפה המשותפת.

כיוון מחקר רביעי הוא שיפור יכולת ההכללה של מודל רשת המיזוג באמצעות למידת העברה (\en{transfer learning}) והרחבת מערך הנתונים. איסוף נתונים ממגוון רחב יותר של סביבות גיאוגרפיות, עונות שנה, ותנאי מזג אוויר יאפשר אימון מודל חזק יותר המסוגל לפעול בתנאים מגוונים. כמו כן, שימוש בטכניקות \en{domain adaptation} ו-\en{domain randomization} עשוי לשפר את יכולת ההכללה מבלי לדרוש איסוף נתונים נוסף.

כיוון מחקר חמישי הוא פיתוח גרסה קלת משקל של המערכת לרחפנים קטנים. גרסה זו תתבסס על מצלמה תרמית ברזולוציה נמוכה יותר (160x120) ומודל רשת מיזוג דחוס (\en{compressed}) באמצעות טכניקות \en{pruning}, \en{quantization}, ו-\en{knowledge distillation}. הערכה ראשונית מראה שניתן להקטין את גודל המודל פי 10 תוך אובדן דיוק של 5-8\% בלבד.

לבסוף, אנו מציעים לחקור את השילוב של מידע סמנטי במערכת המיזוג. מודל זיהוי סמנטי (\en{semantic segmentation}) המותאם לתמונות תרמיות יכול לספק מידע על קטגוריית המכשול (הולך רגל, רכב, עץ) והתנהגותו הצפויה. מידע זה יכול לשפר את קבלת ההחלטות במצבים מורכבים, כגון זיהוי הולך רגל העומד לחצות את נתיב הטיסה.

\subsection{סיכום ומבט לעתיד}
\label{sec:ch09_final}

מחקר זה מדגים את הפוטנציאל הרב של מיזוג תרמי-חזותי בהשראת מערכת החוש התרמי של הצפע הגומתי לניווט אוטונומי בתנאי לילה. המערכת שפותחה משיגת ביצועים עדיפים משמעותית על פני שיטות קיימות, תוך עמידה בדרישות ההפעלה בזמן אמת על פלטפורמת רחפן מוגבלת משאבים. תוצאות אלה מתיישבות עם מגמות מחקר עדכניות בתחום המיזוג הרב-מודאלי לרובוטיקה אוטונומית \cite{Thrun2005ProbRobotics, Kalman1960}.

ההשראה הביולוגית ממנגנון קליטת האינפרה-אדום של הצפע הגומתי אפשרה לנו לפתח גישה ייחודית למיזוג חיישנים, המשלבת עקרונות של עיבוד עצבי עם טכניקות למידה עמוקה מתקדמות. גישה זו מהווה דוגמה לכיצד חקר מערכות ביולוגיות יכול להוביל לפריצות דרך טכנולוגיות בתחום הרובוטיקה האוטונומית.

אנו מאמינים כי המשך הפיתוח בכיווני המחקר העתידיים שהוצעו יוביל למערכות ניווט אוטונומיות ליליות חזקות, יעילות ואמינות יותר, שיאפשרו יישומים חדשים בתחומי חיפוש והצלה, אבטחה, ניטור סביבתי, וחקלאות מדויקת.

\appendix
\section{רשימת סמלים ומשתנים}
\label{sec:appendix_symbols}

\begin{table*}[htbp]
\centering
\caption{רשימת סמלים ומשתנים נפוצים במאמר זה.}
\label{tab:appendix_symbols}
\adjustbox{max width=\columnwidth}{%
\begin{tabular}{@{}lll@{}}
\toprule
\textbf{סמל} & \textbf{הגדרה} & \textbf{יחידות} \\
\midrule
$\mathbf{x}_k$ & וקטור המצב בזמן $k$ & משתנה \\
$\mathbf{z}_k$ & וקטור המדידות בזמן $k$ & משתנה \\
$\mathbf{u}_k$ & וקטור הכניסות (בקרה) בזמן $k$ & משתנה \\
$\mathbf{K}_k$ & רווח קלמן בזמן $k$ & חסרת ממד \\
$\mathbf{P}_k$ & מטריצת השונות המשותפת של השגיאה בזמן $k$ & תלויה במשתנים \\
$\mathbf{R}_k$ & מטריצת השונות המשותפת של רעש המדידה & תלויה במשתנים \\
$\mathbf{Q}_k$ & מטריצת השונות המשותפת של רעש התהליך & תלויה במשתנים \\
$\mathbf{H}_k$ & יעקוביאן מודל המדידה & חסרת ממד \\
$\mathbf{G}_k$ & יעקוביאן מודל התנועה & חסרת ממד \\
$\pi(\cdot)$ & פונקציית היטל מצלמה & פיקסלים \\
$\mathbf{T}_{ab}$ & טרנספורמציה קשיחה ממערכת $a$ ל-$b$ & מטרים, רדיאנים \\
$\mathbf{R}_{ab}$ & מטריצת סיבוב ממערכת $a$ ל-$b$ & חסרת ממד \\
$\mathbf{t}_{ab}$ & וקטור תרגום ממערכת $a$ ל-$b$ & מטרים \\
$\alpha(t)$ & משקל מיזוג אדפטיבי בזמן $t$ & חסרת ממד (0-1) \\
$L_{\text{amb}}(t)$ & עוצמת תאורה סביבתית בזמן $t$ & \en{lux} \\
$T(\mathbf{p})$ & טמפרטורה בפיקסל $\mathbf{p}$ & $^{\circ}$C \\
$S_{\text{therm}}(\mathbf{p})$ & מפת חום תרמית בפיקסל $\mathbf{p}$ & חסרת ממד \\
$C(\mathbf{p})$ & מפת עלות לניווט בפיקסל $\mathbf{p}$ & חסרת ממד \\
$\text{ATE}$ & שגיאת מסלול מוחלטת & מטרים \\
$\text{RPE}$ & שגיאת יחס יחסי & מטרים \\
$\text{PSNR}$ & יחס אות לרעש שיא & \en{dB} \\
$\text{SSIM}$ & מדד דמיון מבני & חסרת ממד (0-1) \\
$\mathcal{N}(\mu, \Sigma)$ & התפלגות נורמלית רב-משתנית & : \\
$\rho(\cdot)$ & פונקציית אובדן חסינה (\en{Huber}) & : \\
$\boxplus$ & אופרטור קומפוזיציה על יריעה & : \\
$\mathfrak{se}(3)$ & אלגברת לי של $SE(3)$ & : \\
\bottomrule
\end{tabular}%
}
\end{table*}